%% ma: L12-B08-0154
%% lop: 12
%% chuong: 2
%% bai: 8
%% dang: 12.8.51
%% loai: TN4
%% muc: TH
%% dap_an: "D"
%% nguon: "CĐ2.4.pdf giáo viên gửi 27/10/2026 (M07, Lê Minh Khôi – Nguyễn Anh Khôi, Chương II Chủ đề 4 Các bài toán về tâm tỉ cự), Câu 4"
%% kien_thuc: "Bài 8 (biểu thức tọa độ của phép toán vectơ, độ dài vectơ), Bài 6 (điểm thỏa hệ thức vectơ, quy tắc ba điểm)"
%% trang_thai: cho-duyet
%% ly_do: "Dạng mới 12.8.51 đề xuất, chờ giáo viên duyệt"
\begin{ex}
Trong không gian $Oxyz$, cho bốn điểm $A(2;-3;7)$, $B(0;4;1)$, $C(3;0;5)$ và $D(3;3;3)$. Tọa độ của điểm $M$ thỏa mãn $\left|\overrightarrow{MA}+\overrightarrow{MB}+\overrightarrow{MC}+\overrightarrow{MD}\right|$ đạt giá trị nhỏ nhất là
\choice{$M(0;1;-4)$}{$M(2;1;0)$}{$M(0;1;-2)$}{\True $M(2;1;4)$}
\loigiai{Gọi $G$ là trọng tâm tứ diện $ABCD$: $G\left(\dfrac84;\dfrac44;\dfrac{16}4\right)=(2;1;4)$.

$\left|\overrightarrow{MA}+\overrightarrow{MB}+\overrightarrow{MC}+\overrightarrow{MD}\right|=4MG$ nhỏ nhất khi $M\equiv G(2;1;4)$.}
\end{ex}
